\SourceLeaf{347}

\begin{center}
[II.]

\textbf{THEOREMS AND PROBLEMS}

\textbf{CONCERNING COMPLEX NUMBERS.}
\end{center}

\Needspace{7\baselineskip}1. $\dfrac{a}{b-\rho}$ represents a circle when $a$, $b$ are fixed numbers and $\rho$ denotes
a variable cyclic number, with
\[
\dfrac{ab}{bb-\beta\beta}
\]
as the centre when $\beta$ is the cyclic factor contained in $b$, and
\[
\dfrac{a}{bb-\beta\beta}
\]
(or the modulus of this quantity) as the radius of the circle.

\Needspace{7\baselineskip}2. (before 1.) If, as above, $\beta$ is the cyclic factor contained in $b$, then
\[
b + \beta\rho \quad \text{and} \quad b + \dfrac{\beta}{\rho}
\]
have the same modulus, and likewise
\[
b + \beta\rho \quad \text{and} \quad \dfrac{\rho\beta\beta}{b\rho+\beta}
\]
have the same clise.

\Needspace{7\baselineskip}3. $\rho - \beta$ and $i\sqrt{\rho\beta}$ have the same clise.

\Needspace{7\baselineskip}4. $\Omega = \dfrac{a\rho+b}{1-\rho}$ represents a straight line. We have
\[
2\Omega = \dfrac{(a+b)(1+\rho) - (a-b)(1-\rho)}{1-\rho}
\]
\[
= (a+b)t - (a-b),
\]
if we set
\[
t = \dfrac{1+\rho}{1-\rho}
\]
so that $t$ becomes a variable imaginary number. We further have
\[
\dfrac{2\Omega}{a+b} = t - \dfrac{a-b}{a+b}.
\]

The value of this quantity whose modulus is a minimum is evidently
just the real part of $\dfrac{a-b}{a+b}$.

Thus if $k$ is the real part of $\dfrac{a-b}{a+b}$, the
sought value of $\Omega$ is
\[
= \tfrac{1}{2} k (a+b).
\]
